\documentclass[11pt,a4paper]{article}
\usepackage[T1]{fontenc}
\usepackage{lmodern,amsmath,amssymb,amsthm,microtype,booktabs}
\usepackage[margin=25mm]{geometry}
\usepackage[colorlinks=true,urlcolor=blue,linkcolor=blue]{hyperref}
\newtheorem{proposition}{Proposition}
\newcommand{\E}{\mathcal E}
\newcommand{\dd}{\mathop{}\!\mathrm d}
\newcommand{\B}{\mathfrak B}
\title{Chiral pairing in the pure-recoil graph\\and the limits of a single-action inverse transform}
\author{Alan Day \and Chengheng\thanks{AI research assistant (OpenAI).}}
\date{Working research note --- 31 October 2026}
\begin{document}\maketitle
\begin{abstract}
We combine the previously controlled internal Barnes graph with
the analytically continued vertex factors and its opposite chiral
partner. Equal probe masses make the internal integral even in the
energy transfer. On a specified imaginary-energy slice this implies
an internal modulus square, but not a new modulus square for every
factor. The complete pure-recoil topology graph retains both Bessel
exponents. A finite energy window already obstructs identifying its
unweighted imaginary-axis transform with the transform of a single
reduced geometric action. This is a self-reviewed model calculation;
the physical external kernel and its normalization remain unspecified.
\end{abstract}

\section{Source-level distinctions and the external interface}
We use the fixed-mass \(\kappa\to\infty\) limit of the CLP
tetrahedron-complement proposal \cite{clp}, then take large masses.
The radius is one. No finite-\(G\) dynamical four-point function is
identified here.
Complex masses refer to continuation of the already obtained recoil-limit
graph; no interchange with a growing-mass finite-\(\kappa\) limit is asserted.
CLP (3.35)--(3.37) distinguish the physical two-point insertion
from the complement partition function. Their normalization
matches a fixed-angle, fixed-probe heavy-heavy-light limit.
It does not specify a full four-point insertion rule or a uniform
limit as the angle approaches \(2\pi\).

CSTY (2.1)--(2.3) \cite{csty} distinguish an energy-domain kernel
from a mixed representation at fixed average observer mass.
With
\[
 m_{1,3}=M+\omega/2\pm u,\qquad m_{2,4}=M-\omega/2\pm v,
\]
the absolute Jacobian is four. The antipodal angles in (2.73) give
\(\sum_i(\theta_i-\theta_{i-1})m_i=2\pi M-2it\omega\), hence
\[
 \mathcal F(t|M)=\frac4{Z_{\rm sphere}(M)}
 \int \dd\omega\,e^{-2it\omega}\int_{\mathcal D_\omega}\dd u\,\dd v\,
       \boldsymbol{\mathcal F}(m_i).                              \tag{1}
\]
The physical kernel, spectral domain and transported cycles are
additional inputs; the complement graph does not determine them alone.
The inverse Laplace contours for lengths \(\beta,\beta_i\)
do not themselves prescribe imaginary contours for the masses.

CLP (3.81), including footnote 12, specifies an inverse transform
of a \emph{reduced semiclassical action}. Identifying it with the full
finite-\(\kappa\) amplitude needs further contour information.
The two time contours in CSTY (2.82)--(2.83) are yet another
operation. These distinctions are part of the original papers.

\newpage
\section{Schwarz pairing and the complete topology ratio}
For a germ continued from real masses, define
\(f^\sharp(z)=\overline{f(\bar z)}\).
The two chiral factors combine as \(A_+(m)A_+^\sharp(m)\).
A fresh \(|A_+(m)|^2\) at complex \(m\) is generally incorrect.
Branches and integration contours are transported together.

For equal probes, the four numerator mass parameters in the internal
Barnes representation are
\(a_j=M+\nu/2+(\pm u\pm v)/2\).
The three right parameters are \(M+\nu\pm\omega/2,2M\).
The first set is independent of \(\omega\); the latter is invariant
under its sign reversal. On a common continued domain,
\[
 \B_+(\omega,u,v)=\B_+(-\omega,u,v).
\]
The same parameter symmetry holds in the original internal
double-sine expression; the numerical checks here concern its
pure-recoil limit only.
For real \(M,\nu,u,v\) and \(\omega=i\alpha\), it follows that
\[
 \B_+(i\alpha,u,v)\B_+^\sharp(i\alpha,u,v)
       =|\B_+(i\alpha,u,v)|^2.                                    \tag{2}
\]
This is a consequence of Schwarz pairing and evenness,
not a general rule for complex external labels.

The previously derived combined vertex \(T\), continued from positive
real masses, gives
\begin{align*}
 \mathcal V(\omega,u,v)
 &=\frac{\prod_iT(m_i,m_{i+1},\nu)}{T(M,M,\nu)^4}
   =P\exp\!\left(\Delta\sum_i\log\E_i\right),\\
 P&=\frac{[(M+\omega/2)^2-u^2][(M-\omega/2)^2-v^2]}{M^4}.
\end{align*}
Here \(\E_i\) is the exact hyperbolic correction written with
\(\log(1+e^{-\pi\ell})\) and \(\log(1-e^{-2\pi m})\), on a strict
positive-real-part triangle-gap domain.
We define the normalized \emph{topology graph}
\[
 \widehat Z(\omega,u,v)=\mathcal V(\omega,u,v)
 \frac{\B_+(\omega,u,v)\B_+^\sharp(\omega,u,v)}
 {|\B_+(0,0,0)|^2}.                                               \tag{3}
\]
Only mass-independent recoil normalization constants cancel here.
This is not normalization by a disconnected physical four-point function.
No vertex is counted twice: \(\B\) is the internal integral.

Writing \(\B=C(u,v)I_{u,v}(\alpha)\), (2) gives
\[
 \widehat Z(i\alpha,u,v)=\mathcal V(i\alpha,u,v)
 \left|\frac{C(u,v)I_{u,v}(\alpha)}{C(0,0)I_{0,0}(0)}\right|^2 .
\]
In general \(\mathcal V(i\alpha,u,v)\) is complex.
For \(M=20,\nu=5,u=0.6,v=1.1,\omega=1.3i\), its value is
\(0.998156494-0.000138120i\).
Taking a fresh modulus square of the continued chiral vertex instead
gives \(1.000916755\).
The correct expression was checked independently against
the Gamma representation of both chiral vertices.

\newpage
\section{Two exponents on the symmetric imaginary axis}
Set \(u=v=0\), \(K=\nu^2/(2M)\), \(\rho=K^3/\nu^2\), and \(\alpha=Ky\).
Take \(K\to\infty,\rho\to0\), with bounded \(y\ge0\).
The vertex ratio is positive and tends to one; \(C(0,0)\) is independent
of \(\alpha\). Thus
\[
 \widehat Z(i\alpha,0,0)=\mathcal V(i\alpha,0,0)|I(\alpha)/I(0)|^2 .
\]
The previous joint estimate reads \(I=e^{iS/\nu}J+e\), \(|e|\le\delta\E\),
with \(\delta=O(\rho+e^{-cK})\).
It gives the precise squared-error inequality
\[
 \big||I|^2-|J|^2\big|\le2\delta\E|J|+\delta^2\E^2.                 \tag{4}
\]
At a zero of \(J\) the bound is \(O(\delta^2\E^2)\); it must not
be divided by \(J\).
Since \(|J(K,0)|\ge\E(K,0)/\sqrt2\), the reference denominator
remains controlled.

Let \(A_K=KK_\alpha(K)\), \(A_I=KI_\alpha(K)\).
The exact Bessel pair is
\[
 |J|^2=A_K^2+\pi^2A_I^2+
                   2\pi\sin(\pi\alpha)A_KA_I.                    \tag{5}
\]
For \(F(y)=\sqrt{1+y^2}-y\operatorname{arsinh}y\),
standard Bessel expansions \cite{dlmf} and (4) imply the
envelope-absolute formula
\begin{align*}
 \widehat Z(iKy,0,0)
 &=\frac{e^{-2K}}{\sqrt{1+y^2}}
   \left[e^{2KF(y)}+e^{-2KF(y)}+2\sin(\pi Ky)\right]\\
 &\quad+O_Y\!\left((\rho+K^{-1})e^{-2K}
                         [e^{KF(y)}+e^{-KF(y)}]^2\right).          \tag{6}
\end{align*}
Near cancellation, use (4)--(5) rather than a relative Debye expansion.
For fixed \(y\ne y_c\), where \(F(y_c)=0\),
\[
 \frac{1}{2K}\log\widehat Z(iKy,0,0)\longrightarrow |F(y)|-1 .
                                                                    \tag{7}
\]
In contrast, the reduced geometric action is
\(W(iKy)=2K[F(y)-1]\).
It corresponds to the first exponential in (6).
For \(y>y_c\), the second exponential dominates; their ratio grows
as \(e^{-4KF(y)}\), not as a mere prefactor.

Numerically \(y_c=1.50887956154\).
The reduced saddle \(y_*=\sinh t\) crosses it at
\(t_c=1.19967864026\).
This is an auxiliary branch-comparison scale, not a physical critical
time or a proof of validity below \(t_c\).
The complete graph divided by
\(e^{2K[F(y)-1]}/\sqrt{1+y^2}\) gives:
\begin{center}
\begin{tabular}{rrr}
\toprule
\(K\)&\(y=1\)&\(y=2\)\\
\midrule
8&0.976384&\(1.04233\times10^9\)\\
16&0.990311&\(1.20309\times10^{18}\)\\
32&0.996309&\(1.53988\times10^{36}\)\\
64&0.998893&\(2.46605\times10^{72}\)\\
\bottomrule
\end{tabular}
\end{center}

\newpage
\section{A finite-window obstruction to direct substitution}
No unbounded-energy extrapolation is needed for the following statement.
Fix \(y_c<a<b\) and \(t\ge0\).
Use the \emph{complete} pure-recoil graph in
\[
 \mathcal W_{a,b}(t)=\int_a^b\dd y\,e^{2Kty}\widehat Z(iKy,0,0).
\]
On this compact window, the second branch dominates uniformly.
Define
\[
 H_t(y)=2[-F(y)-1+ty],\qquad
 H'_t(y)=2[\operatorname{arsinh}y+t]>0 .
\]
The relative graph error is \(O(\rho)\) on the window.
The ordinary endpoint Laplace estimate therefore gives:
\begin{proposition}
In the joint regime above,
\[
 \mathcal W_{a,b}(t)=
 \frac{e^{KH_t(b)}}{2K(\operatorname{arsinh}b+t)\sqrt{1+b^2}}
 \left[1+O_{a,b,t}(\rho+K^{-1})\right].                            \tag{8}
\]
\end{proposition}
The measure is \(\dd y\). Replacing it by \(\dd\alpha\) multiplies
the answer by \(K\).
For \(a=2.5,b=3\),
\[
 H_0(3)=2.58612343506,\qquad H_1(3)=8.58612343506.
\]
These exceed the reduced-action full-axis saddle rate
\(2(\cosh t-1)\), respectively zero and approximately \(1.08616\).
Consequently the unweighted complete topology graph on this same
positive imaginary-axis prescription cannot simply inherit the
single-action result.
This does not prove divergence of the finite-mass exact graph integral,
nor identify a physical external correlator.

A separate, simpler fact concerns only the auxiliary Bessel model
\(P_K(\alpha)=|J(K,\alpha)/J(K,0)|^2\).
At fixed \(K>0\), large positive order gives
\[
 K_\alpha(K)\sim\tfrac12\Gamma(\alpha)(K/2)^{-\alpha},\qquad
 I_\alpha(K)\sim(K/2)^\alpha/\Gamma(\alpha+1).
\]
Thus \(\int_{\mathbb R}\dd\alpha\,e^{2t\alpha}P_K(\alpha)\)
diverges for every real \(t\).
The bounded-\(\alpha/K\) joint estimate cannot be used to transfer
this statement to the finite-\(\nu\) exact graph.
It does prohibit integrating the complete auxiliary two-branch formula
to infinity as a substitute for that graph.

We directly integrated the auxiliary pair on \([2.5,3]\).
Its integral divided by the endpoint leading expression in (8) is:
\begin{center}
\begin{tabular}{rrr}
\toprule
\(K\)&\(t=0\)&\(t=1\)\\
\midrule
8&0.988159&0.980810\\
16&0.994575&0.990951\\
32&0.997395&0.995596\\
\bottomrule
\end{tabular}
\end{center}
The exact-graph statement (8) follows from uniform analysis;
this table is not a new exact Barnes multidimensional quadrature.

\newpage
\section{Transverse consequence and physical calibration}
The previous transverse result used real \(u,v\) with
\(|u|+|v|\le\theta\nu\), fixed \(\theta<1\).
For fixed \(y>y_c\), define
\[
 q=\frac{k}{K}
 =\sqrt{[1-(u+v)^2/\nu^2][1-(u-v)^2/\nu^2]},\quad
 \phi_y(q)=\sqrt{q^2+y^2}-y\operatorname{arsinh}(y/q),
\]
and \(\Delta=\phi_y(1)-\phi_y(q)>0\) for nonzero transverse ratios.
The analytic vertex ratio tends to one in the joint regime.
Schwarz pairing doubles the internal growth exponent:
\[
 \left|\frac{\widehat Z(iKy,u,v)}{\widehat Z(iKy,0,0)}\right|
 =\left(\frac{1+y^2}{q^2+y^2}\right)^{1/2}e^{2K\Delta}
       [1+O_{\theta,y}(\rho+K^{-1})].                             \tag{9}
\]
For \(y=2,u=0.3\nu,v=0.2\nu\), \(2\Delta=0.65667292183\).
The computed complete graph ratios have magnitudes
approximately \(192.703\), \(37040.448\), and \(1.359471\times10^9\)
at \(K=8,16,32\).
Their complex values are preserved: the last has imaginary part
approximately \(8.10344\), which is small relatively but is not zero.
Equation (9) is not physical time growth.

\subsection*{What is now determined}
The two chiral factors of the specified topology graph can be paired
without discarding a branch. The exact vertex continuation is available
in the recoil limit, and the joint estimates control its internal
energy slices. These results strengthen the earlier conclusion that
a central real transverse Gaussian cannot describe the entire graph.

\subsection*{What is still missing}
The physical insertion legs, observer spectral weights and disconnected
normalization must convert a complement amplitude into the
energy-domain kernel in (1).
The trace interpretation itself is part of the physics under
investigation. No physical state, algebra, positivity theorem or
finite-\(G\) microscopic model is newly established by (2)--(9).
The original papers' contour qualifications remain applicable.

The next concrete calibration is analytic pole structure.
Starting from the original Barnes integral, locate its
energy-transfer pinch poles and check what happens when the four
combined vertices and the opposite chirality are included.
Compare the result with the direct \(C_0\) factor in the known physical
two-point prescription and with the disconnected limit.
This can test whether a proposed external-leg conversion is necessary,
without assuming that the complement amplitude already is a physical
four-point kernel.

If a sewing or insertion rule is proposed, list the extra spectral
measure and inner-product assumptions explicitly. Verify the
two-point calibration and a factorization limit before using that
rule to choose an external energy cycle. Rotating \(u,v\) merely to
make a growing branch decay would not supply this missing derivation.

\newpage
\section{Evidence and reproducibility}
The final program records 37 checks:
two symbolic external-variable identities; six original Barnes-kernel
evenness checks; six Gamma-vertex or Schwarz identities;
four noninteger-order Bessel interference identities;
twelve paired-vertex precision comparisons; one crossover equation;
and six window-quadrature refinements.
The general algebra uses 90 digits, with vertex checks at 110 digits.
The window computation uses orders 20 and 32 at 55 and 75 digits.

Twelve symmetric and three transverse graph ratios reuse previously
frozen complex branch evaluations. This is new composition with the
analytic vertex factors, not a claim to have repeated exact full
quadrature at all those points.
At an earlier auxiliary cancellation sample,
the recorded paired internal magnitude is
\[
 |I|^2/\E^2\simeq9.24402\times10^{-5}.
\]
No division by a vanishing leading term is used.

The first 37-check run failed four window refinement checks.
A single quadrature panel did not resolve the sharp endpoint
well enough: the largest relative change was
\(3.05825\times10^{-7}\), above the \(10^{-20}\) criterion.
The initial code and JSON are retained.
Subdividing to maximum panel width \(0.1\), without relaxing the
criterion, gives a maximum relative change \(1.25050\times10^{-24}\).
These are empirical diagnostics, not interval certificates.

The complete-graph window proposition is a consequence of the earlier
uniform internal theorem, the explicit vertex formula, and endpoint
Laplace analysis. The auxiliary whole-axis divergence uses a distinct
fixed-\(K\), large-order limit. These two claims should not be conflated.

\small
Project-relative records:
\begin{quote}\footnotesize
code/recoil\_pairing.py\\
results/stage4/recoil-pairing/diagnostics.json\\
results/stage4/recoil-pairing/initial-window-grid.json\\
results/stage4/recoil-pairing/initial-window-grid.py\\
notes/stage4-recoil-pairing.md\\
literature/stage4-recoil-pairing-sources.md
\end{quote}
\begin{thebibliography}{9}
\bibitem{clp} J. Chandra, D. Liska and B. Post,
\emph{3D de Sitter from complex Virasoro TQFT: Observers and knots},
\href{https://arxiv.org/abs/2610.30247v1}{arXiv:2610.30247v1}.
See (3.35)--(3.37), (3.43)--(3.46), (3.77)--(3.84), and footnote 12.
\bibitem{csty} Y. Chen, D. Stanford, H. Tang and Z. Yang,
\emph{Negative shocks versus static patch holography},
\href{https://arxiv.org/abs/2607.14042v1}{arXiv:2607.14042v1}.
See Section 2.1, (2.73), (2.82)--(2.83), footnote 19,
and the external-contour discussion in Section 2.6.
\bibitem{dlmf} NIST DLMF,
\href{https://dlmf.nist.gov/10.40}{Section 10.40},
\href{https://dlmf.nist.gov/10.41}{Section 10.41}.
The standard Bessel asymptotics are not new results.
\end{thebibliography}
\end{document}
